\BourbakiFrontChapter{引言}

数学的多数分支都涉及一类结构，它们不同于本丛书《代数》一书所研究的代数结构（群、环、域等）：即给极限、连续与邻域这些直观概念以数学内容的结构。这些结构正是本书的研究对象。

从历史上看，极限与连续的思想在数学中出现得很早，尤其是在几何学中；随着分析学的发展及其在实验科学中的应用，它们的作用不断增长，因为这些思想与实验测定和近似的思想密切相关。但由于大多数实验测定都是度量，亦即对一个或多个数的测定，因此不足为怪的是，数学中的极限与连续概念起初只出现在实数理论及其衍生物与应用领域（复数、实变量或复变量的实值或复值函数、欧氏几何及有关几何）之中。

近来人们认识到，这些思想的适用范围远远超出了经典分析中的实数与复数（见第一章的历史注记）。通过分析与抽象的努力，它们的本质内容已被提炼出来，其结果是得到一个工具，它在数学的许多分支中的用处已经显而易见。

为了揭示极限、连续与邻域这些思想中的本质内容，我们将从分析邻域概念入手（尽管从历史上看，它比另外两个概念出现得更晚）。如果我们从物理上的近似概念出发，自然可以说：集合 E 的子集 A 是 A 中元素 a 的邻域，如果每当我们用某个“近似”a 的元素代替 a 时，这个新元素也属于 A，当然要假定所涉及的“误差”足够小；或者换句话说，如果 E 中所有“充分接近”a 的点都属于 A。每当能够给“充分小的误差”概念或“一个元素充分接近另一个元素”概念以精确意义时，这个定义就是有意义的。沿这个方向，第一个想法是假定两个元素之间的“距离”可以用一个（正）实数来度量。一旦定义了集合中任意两个元素之间的“距离”，应当如何定义元素 a 的“邻域”便是清楚的了：一个子集将是 a 的邻域，如果它包含所有与 a 的距离小于某个事先给定的严格正数的元素。当然，除非对“距离”附加某些条件或公理（例如，欧氏几何中关于三角形三个顶点之间的距离所成立的不等式，对我们的推广距离仍应成立），否则不能指望由这个定义发展出一个有趣的理论。这样我们就得到了欧氏几何的一个巨大推广。继续使用几何学的语言是方便的：于是，定义了“距离”的集合，其元素称为点，而该集合本身称为空间。我们将在第九章研究此类空间。

到目前为止，我们尚未成功摆脱实数。然而，如此定义的空间具有许许多多性质，陈述这些性质无须涉及产生它们的那个“距离”。例如，包含 a 的一个邻域的每个子集仍是 a 的邻域，a 的两个邻域的交是 a 的邻域。这些性质以及其他性质有大量推论，无须再求助于当初使我们得以定义邻域的“距离”就能推导出来。我们得到的是其中不再提及量或距离的命题。

这样我们终于被引到拓扑空间的一般概念，它不依赖于任何预备的实数理论。我们说集合 E 赋有一个拓扑结构，只要我们以某种方式使 E 的每个元素都对应了 E 的一个子集族，族中的子集称为该元素的邻域——当然要假定这些邻域满足某些条件（拓扑结构的公理）。显然，所加公理的选择在某种程度上是任意的，并且在历史上曾是大量试验的主题（见第一章的历史注记）。最终得到的公理体系对于数学当前的需要而言足够宽泛，又不至于陷入过分而无益的一般性。

赋有拓扑结构的集合称为拓扑空间，其元素称为点。研究拓扑结构的数学分支名称是拓扑学（按词源是“位置的学科”，这并不是一个特别有表现力的名称），如今人们更愿意用这个名称，而不用较早的同义名称 Analysis situs（位置分析）。

为了表述邻域的思想，我们曾从“一个元素充分接近”另一个元素这一模糊概念出发。反过来，拓扑结构现在使我们能够给“某某性质对所有充分接近 a 的点成立”这一说法以精确的意义：按定义，这意味着具有该性质的点所成的集合，在所论的拓扑结构下是 a 的一个邻域。

从邻域概念引出一系列其他概念，对它们的研究属于拓扑学：集合的内部、集合的闭包、集合的边界、开集、闭集等等（见第一章 § 1）。例如，子集 A 是开集，如果只要点 \(a\) 属于 A，所有充分接近 \(a\) 的点就都属于 A；换句话说，如果 A 是它的每个点的邻域。邻域的公理对所有这些概念都有某些推论；例如，两个开集的交是开集（因为我们已假定了 \(a\) 的两个邻域的交是 \(a\) 的邻域）。反过来，我们也可以从这些导出概念之一出发，而不从邻域概念出发；例如，可以假定开集是已知的，并把开集族的性质取作公理（其中一条性质刚才已作为例子陈述过）。然后可以验证，由开集的知识能够重新构造出邻域；邻域的公理现在成为我们取作出发点的新的开集公理的推论。于是，拓扑结构可以用种种基本等价的不同方式来定义。本书将从开集概念出发，因为相应的公理最简单。

一旦定义了拓扑结构，就容易使连续性的观念精确化。直观上，如果只要变元保持在所论点的充分近处，函数的值就变化得要多小有多小，则函数在该点连续。于是，只要函数的变元空间与值空间都是拓扑空间，连续性就有精确的意义。精确的定义在第一章 § 2 中给出。

与连续性一样，极限的观念涉及两个集合，每个都赋有适当的结构，还有一个集合到另一个集合的映射。例如，实数序列 \(a_{n}\) 的极限涉及自然数集合 N、实数集合 R，以及前一集合到后一集合的映射。称实数 a 为该序列的极限，如果无论取 a 的哪个邻域 V，该邻域都包含除 n 的有限多个值以外的所有 \(a_{n}\)；也就是说，如果使 \(a_{n}\) 属于 V 的那些自然数 n 所成的集合是 N 的余集有限的子集。注意，由于谈论的是邻域，R 被假定赋有拓扑结构；至于集合 N，我们让某个特定的子集族扮演特殊角色，即余集有限的那些子集。这是一个一般事实：每当我们谈论极限，考虑的都是映射 \(f\)（从一个集合 \(\mathbf{E}\) 到一个拓扑空间 \(\mathbf{F}\)），并且我们说 \(f\) 以点 \(a\)（\(\mathbf{F}\) 中的点）为极限，如果由这样的元素 \(x\)（属于 \(\mathbf{E}\)）所组成的集合——其像 \(f(x)\) 属于

a 的某个邻域 V {[}此集合正是“逆像” \(\overline{f}^{1}(V)\) {]}——无论邻域 V 如何，都属于 E 的某个事先给定的子集族 F。为使极限概念具有通常归之于它的那些本质性质，族 F 必须满足某些公理，这些公理在第一章 § 6 中陈述。E 的这样的子集族 F 称为 E 上的滤子。滤子的概念与极限的概念不可分离，它也出现在拓扑学的其他场合；例如，拓扑空间中一点的诸邻域构成一个滤子。

对所有这些概念的一般研究是第一章的主要目的。此外，那里还考虑拓扑空间的一些特殊类：满足更多限制性公理的空间，或由其他给定空间通过特定程序得到的空间。

如前所述，集合上的拓扑结构使人能够给“每当 x 充分接近 a 时，x 就具有性质 \(P\{x\}\)”这一说法以精确的意义。但是，除去已定义“距离”的情形，尚不清楚应当给“每一对彼此充分接近的点 x, y 都具有性质 \(P\{x, y\}\)”这一说法以什么意义，因为先验地我们没有比较两个不同点的邻域的手段。而彼此接近的点对这一概念在经典分析中频繁出现（例如，在涉及一致连续性的命题中）。因此，重要的是我们能够在完全一般的情形下给这一概念以精确意义，于是我们被引导去定义比拓扑结构更丰富的结构，即一致结构。它们是第二章的主题。

本卷的其余各章致力于这样的问题：在这些问题中，除拓扑结构或一致结构外，还出现别的结构。例如，赋有适当拓扑（在某种意义上与群结构相容）的群称为拓扑群。拓扑群在第三章中研究，我们将在那里特别看到，每个拓扑群如何可以赋有某些一致结构。

在第四章中，我们把上述原理应用于有理数域。这使我们能够定义实数域；由于其重要性，我们将相当详细地研究它。在随后各章中，我们从实数出发定义某些拓扑空间，它们在拓扑学对经典几何的应用中具有特殊意义：有限维向量空间、球面、射影空间等。我们还考虑与

实数的加法群密切相关的某些拓扑群，我们对它们作公理化的刻画，而这就把我们引导到经典分析中最重要的函数（指数函数、对数函数与三角函数）的定义及初等性质。

在第九章中我们回到一般拓扑空间，但现在有一种新工具，即实数，可供我们支配。我们特别研究拓扑借助于“距离”来定义的空间；这些空间具有一些性质，其中若干不能推广到更一般的空间。在第十章中我们研究拓扑空间到一致空间的映射所成的集合（函数空间）；这些集合在赋以适当的拓扑后具有有趣的性质，这些性质在经典分析中已经起着重要作用。

